
% https://akaedu.github.io/book/ch33.html

\chapter{信号}
\BookAnchor{db.ch33}\BookAnchor{db.ch33.id2883851}%


\BookNav{\hyperlink{db.ch32s05}{上一页}\enspace\textperiodcentered\enspace \hyperlink{db.pt03}{上一级}\enspace\textperiodcentered\enspace \hyperlink{db.index}{起始页}\enspace\textperiodcentered\enspace \hyperlink{db.ch33s01}{下一页}}


% https://akaedu.github.io/book/ch33s01.html

\section{信号的基本概念}
\BookAnchor{db.ch33s01}\BookAnchor{db.ch33s01.id2883230}%
为了理解信号，先从我们最熟悉的场景说起：


\begin{enumerate}

\item 用户输入命令，在Shell下启动一个前台进程。

\item 用户按下Ctrl-C，这个键盘输入产生一个硬件中断。

\item 如果CPU当前正在执行这个进程的代码，则该进程的用户空间代码暂停执行，CPU从用户态切换到内核态处理硬件中断。

\item 终端驱动程序将Ctrl-C解释成一个\BookCode{SIGINT}信号，记在该进程的PCB中（也可以说发送了一个\BookCode{SIGINT}信号给该进程）。

\item 当某个时刻要从内核返回到该进程的用户空间代码继续执行之前，首先处理PCB中记录的信号，发现有一个\BookCode{SIGINT}信号待处理，而这个信号的默认处理动作是终止进程，所以直接终止进程而不再返回它的用户空间代码执行。
\end{enumerate}


注意，Ctrl-C产生的信号只能发给前台进程。在\hyperlink{db.ch30s03.id2867242}{第 3.\BookSplit{}3 节 “wait和waitpid函数”}中我们看到一个命令后面加个\BookCode{\&}可以放到后台运行，这样Shell不必等待进程结束就可以接受新的命令，启动新的进程。Shell可以同时运行一个前台进程和任意多个后台进程，只有前台进程才能接到像Ctrl-C这种控制键产生的信号。前台进程在运行过程中用户随时可能按下Ctrl-C而产生一个信号，也就是说该进程的用户空间代码执行到任何地方都有可能收到\BookCode{SIGINT}信号而终止，所以信号相对于进程的控制流程来说是异步（Asynchronous）\BookAnchor{db.ch33s01.id2883948}的。



用\BookCode{kill -\BookSplit{}l}命令可以察看系统定义的信号列表：


\BookTerminal{code/ch33s01/01.txt}{9}{11}


每个信号都有一个编号和一个宏定义名称，这些宏定义可以在\BookCode{signal.\BookSplit{}h}中找到，例如其中有定义\BookCode{\#define SIGINT\BookSplit{} 2}。编号34以上的是实时信号，本章只讨论编号34以下的信号，不讨论实时信号。这些信号各自在什么条件下产生，默认的处理动作是什么，在\BookCode{signal(\BookSplit{}7)\BookSplit{}}中都有详细说明：


\BookTerminal{code/ch33s01/02.txt}{9}{8}


上表中第一列是各信号的宏定义名称，第二列是各信号的编号，第三列是默认处理动作，\BookCode{Term}表示终止当前进程，\BookCode{Core}表示终止当前进程并且Core Dump（下一节详细介绍什么是Core Dump），\BookCode{Ign}表示忽略该信号，\BookCode{Stop}表示停止当前进程，\BookCode{Cont}表示继续执行先前停止的进程，表中最后一列是简要介绍，说明什么条件下产生该信号。



产生信号的条件主要有：


\begin{itemize}

\item 用户在终端按下某些键时，终端驱动程序会发送信号给前台进程，例如Ctrl-C产生\BookCode{SIGINT}信号，Ctrl-\textbackslash{}产生\BookCode{SIGQUIT}信号，Ctrl-Z产生\BookCode{SIGTSTP}信号（可使前台进程停止，这个信号将在\hyperlink{db.ch34.id2890980}{第 34 章 \emph{终端、作业控制与守护进程}}详细解释）。

\item 硬件异常产生信号，这些条件由硬件检测到并通知内核，然后内核向当前进程发送适当的信号。例如当前进程执行了除以0的指令，CPU的运算单元会产生异常，内核将这个异常解释为\BookCode{SIGFPE}信号发送给进程。再比如当前进程访问了非法内存地址，，MMU会产生异常，内核将这个异常解释为\BookCode{SIGSEGV}信号发送给进程。

\item 一个进程调用\BookCode{kill(\BookSplit{}2)\BookSplit{}}函数可以发送信号给另一个进程。

\item 可以用\BookCode{kill(\BookSplit{}1)\BookSplit{}}命令发送信号给某个进程，\BookCode{kill(\BookSplit{}1)\BookSplit{}}命令也是调用\BookCode{kill(\BookSplit{}2)\BookSplit{}}函数实现的，如果不明确指定信号则发送\BookCode{SIGTERM}信号，该信号的默认处理动作是终止进程。

\item 当内核检测到某种软件条件发生时也可以通过信号通知进程，例如闹钟超时产生\BookCode{SIGALRM}信号，向读端已关闭的管道写数据时产生\BookCode{SIGPIPE}信号。
\end{itemize}


如果不想按默认动作处理信号，用户程序可以调用\BookCode{sigaction(\BookSplit{}2)\BookSplit{}}函数告诉内核如何处理某种信号（\BookCode{sigaction}函数稍后详细介绍），可选的处理动作有以下三种：


\begin{enumerate}

\item 忽略此信号。

\item 执行该信号的默认处理动作。

\item 提供一个信号处理函数，要求内核在处理该信号时切换到用户态执行这个处理函数，这种方式称为捕捉（Catch）\BookAnchor{db.ch33s01.id2884226}一个信号。
\end{enumerate}

\BookNav{\hyperlink{db.ch33}{上一页}\enspace\textperiodcentered\enspace \hyperlink{db.ch33}{上一级}\enspace\textperiodcentered\enspace \hyperlink{db.index}{起始页}\enspace\textperiodcentered\enspace \hyperlink{db.ch33s02}{下一页}}


% https://akaedu.github.io/book/ch33s02.html

\section{产生信号}
\BookAnchor{db.ch33s02}\BookAnchor{db.ch33s02.id2884237}%
\subsection{通过终端按键产生信号}
\BookAnchor{db.ch33s02.id2884244}%


上一节讲过，\BookCode{SIGINT}的默认处理动作是终止进程，\BookCode{SIGQUIT}的默认处理动作是终止进程并且Core Dump，现在我们来验证一下。



首先解释什么是Core Dump。当一个进程要异常终止时，可以选择把进程的用户空间内存数据全部保存到磁盘上，文件名通常是\BookCode{core}，这叫做Core Dump\BookAnchor{db.ch33s02.id2884282}。进程异常终止通常是因为有Bug，比如非法内存访问导致段错误，事后可以用调试器检查\BookCode{core}文件以查清错误原因，这叫做Post-mortem Debug\BookAnchor{db.ch33s02.id2884298}。一个进程允许产生多大的\BookCode{core}文件取决于进程的Resource Limit（这个信息保存在PCB中）。默认是不允许产生\BookCode{core}文件的，因为\BookCode{core}文件中可能包含用户密码等敏感信息，不安全。在开发调试阶段可以用\BookCode{ulimit}命令改变这个限制，允许产生\BookCode{core}文件。



首先用\BookCode{ulimit}命令改变Shell进程的Resource Limit，允许\BookCode{core}文件最大为1024K：


\BookTerminal{code/ch33s02/01.txt}{3}{1}


然后写一个死循环程序：


\BookSource{numbers=none}{code/ch33s02/02.pv}{9}{7}


前台运行这个程序，然后在终端键入Ctrl-C或Ctrl-\textbackslash{}：


\BookTerminal{code/ch33s02/03.txt}{8}{6}


\BookCode{ulimit}命令改变了Shell进程的Resource Limit，\BookCode{a.\BookSplit{}out}进程的PCB由Shell进程复制而来，所以也具有和Shell进程相同的Resource Limit值，这样就可以产生Core Dump了。


\subsection{调用系统函数向进程发信号}
\BookAnchor{db.ch33s02.id2884400}%


仍以上一节的死循环程序为例，首先在后台执行这个程序，然后用\BookCode{kill}命令给它发\BookCode{SIGSEGV}信号。


\BookTerminal{code/ch33s02/04.txt}{7}{5}


7940是\BookCode{a.\BookSplit{}out}进程的id。之所以要再次回车才显示\BookCode{Segmentation f\BookSplit{}ault}，是因为在7940进程终止掉之前已经回到了Shell提示符等待用户输入下一条命令，Shell不希望\BookCode{Segmentation f\BookSplit{}ault}信息和用户的输入交错在一起，所以等用户输入命令之后才显示。指定某种信号的\BookCode{kill}命令可以有多种写法，上面的命令还可以写成\BookCode{kill -\BookSplit{}SEGV 7940}或\BookCode{kill -\BookSplit{}11 7940}，11是信号\BookCode{SIGSEGV}的编号。以往遇到的段错误都是由非法内存访问产生的，而这个程序本身没错，给它发\BookCode{SIGSEGV}也能产生段错误。



\BookCode{kill}命令是调用\BookCode{kill}函数实现的。\BookCode{kill}函数可以给一个指定的进程发送指定的信号。\BookCode{raise}函数可以给当前进程发送指定的信号（自己给自己发信号）。


\BookSource{numbers=none}{code/ch33s02/05.pv}{6}{4}


这两个函数都是成功返回0，错误返回-1。



\BookCode{abort}函数使当前进程接收到\BookCode{SIGABRT}信号而异常终止。


\BookSource{numbers=none}{code/ch33s02/06.pv}{5}{3}


就像\BookCode{exit}函数一样，\BookCode{abort}函数总是会成功的，所以没有返回值。


\subsection{由软件条件产生信号}
\BookAnchor{db.ch33s02.id2884567}%


\BookCode{SIGPIPE}是一种由软件条件产生的信号，在\hyperlink{db.ch30s04.id2867940}{例 30.\BookSplit{}7 “管道”}中已经介绍过了。本节主要介绍\BookCode{alarm}函数和\BookCode{SIGALRM}信号。


\BookSource{numbers=none}{code/ch33s02/07.pv}{5}{3}


调用\BookCode{alarm}函数可以设定一个闹钟，也就是告诉内核在\BookCode{seconds}秒之后给当前进程发\BookCode{SIGALRM}信号，该信号的默认处理动作是终止当前进程。这个函数的返回值是0或者是以前设定的闹钟时间还余下的秒数。打个比方，某人要小睡一觉，设定闹钟为30分钟之后响，20分钟后被人吵醒了，还想多睡一会儿，于是重新设定闹钟为15分钟之后响，“以前设定的闹钟时间还余下的时间”就是10分钟。如果\BookCode{seconds}值为0，表示取消以前设定的闹钟，函数的返回值仍然是以前设定的闹钟时间还余下的秒数。


\BookExampleTitle{alarm}\BookAnchor{db.ch33s02.id2884661}

\BookSource{numbers=none}{code/ch33s02/08.pv}{9}{11}


这个程序的作用是1秒钟之内不停地数数，1秒钟到了就被\BookCode{SIGALRM}信号终止。

\BookNav{\hyperlink{db.ch33s01}{上一页}\enspace\textperiodcentered\enspace \hyperlink{db.ch33}{上一级}\enspace\textperiodcentered\enspace \hyperlink{db.index}{起始页}\enspace\textperiodcentered\enspace \hyperlink{db.ch33s03}{下一页}}


% https://akaedu.github.io/book/ch33s03.html

\section{阻塞信号}
\BookAnchor{db.ch33s03}\BookAnchor{db.ch33s03.id2884689}%
\subsection{信号在内核中的表示}
\BookAnchor{db.ch33s03.id2884694}%


以上我们讨论了信号\emph{产生}（Generation）的各种原因，而实际执行信号的处理动作称为信号\emph{递达}（Delivery），信号从产生到递达之间的状态，称为信号\emph{未决}（Pending）。进程可以选择阻塞（Block）某个信号。被阻塞的信号产生时将保持在未决状态，直到进程解除对此信号的阻塞，才执行递达的动作。注意，阻塞和忽略是不同的，只要信号被阻塞就不会递达，而忽略是在递达之后可选的一种处理动作。信号在内核中的表示可以看作是这样的：


\BookFigure{信号在内核中的表示示意图}{\BookImage{images/signal.internal.png}{0.832}}{\BookAnchor{db.ch33s03.id2884728}}


每个信号都有两个标志位分别表示阻塞和未决，还有一个函数指针表示处理动作。信号产生时，内核在进程控制块中设置该信号的未决标志，直到信号递达才清除该标志。在上图的例子中，


\begin{enumerate}

\item \BookCode{SIGHUP}信号未阻塞也未产生过，当它递达时执行默认处理动作。

\item \BookCode{SIGINT}信号产生过，但正在被阻塞，所以暂时不能递达。虽然它的处理动作是忽略，但在没有解除阻塞之前不能忽略这个信号，因为进程仍有机会改变处理动作之后再解除阻塞。

\item \BookCode{SIGQUIT}信号未产生过，一旦产生\BookCode{SIGQUIT}信号将被阻塞，它的处理动作是用户自定义函数\BookCode{sighandler}。
\end{enumerate}


如果在进程解除对某信号的阻塞之前这种信号产生过多次，将如何处理？POSIX.\BookSplit{}1允许系统递送该信号一次或多次。Linux是这样实现的：常规信号在递达之前产生多次只计一次，而实时信号在递达之前产生多次可以依次放在一个队列里。本章不讨论实时信号。从上图来看，每个信号只有一个bit的未决标志，非0即1，不记录该信号产生了多少次，阻塞标志也是这样表示的。因此，未决和阻塞标志可以用相同的数据类型\BookCode{sigset\_\BookSplit{}t}来存储，\BookCode{sigset\_\BookSplit{}t}称为信号集，这个类型可以表示每个信号的“有效”或“无效”状态，在阻塞信号集中“有效”和“无效”的含义是该信号是否被阻塞，而在未决信号集中“有效”和“无效”的含义是该信号是否处于未决状态。下一节将详细介绍信号集的各种操作。阻塞信号集也叫做当前进程的信号屏蔽字（Signal Mask）\BookAnchor{db.ch33s03.id2884862}，这里的“屏蔽”应该理解为阻塞而不是忽略。


\subsection{信号集操作函数}
\BookAnchor{db.ch33s03.id2884876}%


\BookCode{sigset\_\BookSplit{}t}类型对于每种信号用一个bit表示“有效”或“无效”状态，至于这个类型内部如何存储这些bit则依赖于系统实现，从使用者的角度是不必关心的，使用者只能调用以下函数来操作\BookCode{sigset\_\BookSplit{}t}变量，而不应该对它的内部数据做任何解释，比如用\BookCode{printf}直接打印\BookCode{sigset\_\BookSplit{}t}变量是没有意义的。


\BookSource{numbers=none}{code/ch33s03/01.pv}{9}{7}


函数\BookCode{sigemptyset}初始化\BookCode{set}所指向的信号集，使其中所有信号的对应bit清零，表示该信号集不包含任何有效信号。函数\BookCode{sigfillset}初始化\BookCode{set}所指向的信号集，使其中所有信号的对应bit置位，表示该信号集的有效信号包括系统支持的所有信号。注意，在使用\BookCode{sigset\_\BookSplit{}t}类型的变量之前，一定要调用\BookCode{sigemptyset}或\BookCode{sigfillset}做初始化，使信号集处于确定的状态。初始化\BookCode{sigset\_\BookSplit{}t}变量之后就可以在调用\BookCode{sigaddset}和\BookCode{sigdelset}在该信号集中添加或删除某种有效信号。这四个函数都是成功返回0，出错返回-1。\BookCode{sigismember}是一个布尔函数，用于判断一个信号集的有效信号中是否包含某种信号，若包含则返回1，不包含则返回0，出错返回-1。


\subsection{sigprocmask}
\BookAnchor{db.ch33s03.id2885022}%


调用函数\BookCode{sigprocmask}可以读取或更改进程的信号屏蔽字。


\BookSource{numbers=none}{code/ch33s03/02.pv}{5}{3}


返回值：若成功则为0，若出错则为-1



如果\BookCode{oset}是非空指针，则读取进程的当前信号屏蔽字通过\BookCode{oset}参数传出。如果\BookCode{set}是非空指针，则更改进程的信号屏蔽字，参数\BookCode{how}指示如何更改。如果\BookCode{oset}和\BookCode{set}都是非空指针，则先将原来的信号屏蔽字备份到\BookCode{oset}里，然后根据\BookCode{set}和\BookCode{how}参数更改信号屏蔽字。假设当前的信号屏蔽字为\BookCode{mask}，下表说明了\BookCode{how}参数的可选值。


\begingroup
\setlength{\tabcolsep}{6pt}
\small
\sloppy
\renewcommand{\arraystretch}{1.16}
\setlength{\parskip}{0pt}
\begin{longtable}{L{\dimexpr 0.182744\linewidth-2\tabcolsep\relax}L{\dimexpr 0.817256\linewidth-2\tabcolsep\relax}}
\BookTableTitle{how参数的含义}\BookTableLabel{\BookAnchor{db.ch33s03.id2885118}}
\toprule
\BookCode{SIG\_\BookSplit{}BLOCK} & \BookCode{set}包含了我们希望添加到当前信号屏蔽字的信号，相当于mask=\BookSplit{}mask|set \\
\addlinespace[2pt]
\BookCode{SIG\_\BookSplit{}UNBLOCK} & \BookCode{set}包含了我们希望从当前信号屏蔽字中解除阻塞的信号，相当于mask=\BookSplit{}mask\&\textasciitilde{}set \\
\addlinespace[2pt]
\BookCode{SIG\_\BookSplit{}SETMASK} & 设置当前信号屏蔽字为\BookCode{set}所指向的值，相当于mask=\BookSplit{}set \\
\bottomrule
\end{longtable}
\endgroup


如果调用\BookCode{sigprocmask}解除了对当前若干个未决信号的阻塞，则在\BookCode{sigprocmask}返回前，至少将其中一个信号递达。


\subsection{sigpending}
\BookAnchor{db.ch33s03.id2885205}%

\BookSource{numbers=none}{code/ch33s03/03.pv}{5}{3}


\BookCode{sigpending}读取当前进程的未决信号集，通过\BookCode{set}参数传出。调用成功则返回0，出错则返回-1。



下面用刚学的几个函数做个实验。程序如下：


\BookSource{numbers=left,xleftmargin=19pt}{code/ch33s03/04.pv}{9}{28}


程序运行时，每秒钟把各信号的未决状态打印一遍，由于我们阻塞了\BookCode{SIGINT}信号，按Ctrl-C将会使\BookCode{SIGINT}信号处于未决状态，按Ctrl-\textbackslash{}仍然可以终止程序，因为\BookCode{SIGQUIT}信号没有阻塞。


\BookTerminal{code/ch33s03/05.txt}{8}{6}

\BookNav{\hyperlink{db.ch33s02}{上一页}\enspace\textperiodcentered\enspace \hyperlink{db.ch33}{上一级}\enspace\textperiodcentered\enspace \hyperlink{db.index}{起始页}\enspace\textperiodcentered\enspace \hyperlink{db.ch33s04}{下一页}}


% https://akaedu.github.io/book/ch33s04.html

\section{捕捉信号}
\BookAnchor{db.ch33s04}\BookAnchor{db.ch33s04.id2885284}%
\subsection{内核如何实现信号的捕捉}
\BookAnchor{db.ch33s04.id2885289}%


如果信号的处理动作是用户自定义函数，在信号递达时就调用这个函数，这称为捕捉信号。由于信号处理函数的代码是在用户空间的，处理过程比较复杂，举例如下：


\begin{enumerate}

\item 用户程序注册了\BookCode{SIGQUIT}信号的处理函数\BookCode{sighandler}。

\item 当前正在执行\BookCode{main}函数，这时发生中断或异常切换到内核态。

\item 在中断处理完毕后要返回用户态的\BookCode{main}函数之前检查到有信号\BookCode{SIGQUIT}递达。

\item 内核决定返回用户态后不是恢复\BookCode{main}函数的上下文继续执行，而是执行\BookCode{sighandler}函数，\BookCode{sighandler}和\BookCode{main}函数使用不同的堆栈空间，它们之间不存在调用和被调用的关系，是两个独立的控制流程。

\item \BookCode{sighandler}函数返回后自动执行特殊的系统调用\BookCode{sigreturn}再次进入内核态。

\item 如果没有新的信号要递达，这次再返回用户态就是恢复\BookCode{main}函数的上下文继续执行了。
\end{enumerate}

\BookFigure{信号的捕捉}{\BookImage{images/signal.catch.png}{0.624}}{\BookAnchor{db.ch33s04.id2885414}}


上图出自\hyperlink{db.bi01.bibli.ulk}{[ULK]}。


\subsection{sigaction}
\BookAnchor{db.ch33s04.id2885439}%

\BookSource{numbers=none}{code/ch33s04/01.pv}{5}{3}


\BookCode{sigaction}函数可以读取和修改与指定信号相关联的处理动作。调用成功则返回0，出错则返回-1。\BookCode{signo}是指定信号的编号。若\BookCode{act}指针非空，则根据\BookCode{act}修改该信号的处理动作。若\BookCode{oact}指针非空，则通过\BookCode{oact}传出该信号原来的处理动作。\BookCode{act}和\BookCode{oact}指向\BookCode{sigaction}结构体：


\BookSource{numbers=none}{code/ch33s04/02.pv}{9}{9}


将\BookCode{sa\_\BookSplit{}handler}赋值为常数\BookCode{SIG\_\BookSplit{}IGN}传给\BookCode{sigaction}表示忽略信号，赋值为常数\BookCode{SIG\_\BookSplit{}DFL}表示执行系统默认动作，赋值为一个函数指针表示用自定义函数捕捉信号，或者说向内核注册了一个信号处理函数，该函数返回值为\BookCode{void}，可以带一个\BookCode{int}参数，通过参数可以得知当前信号的编号，这样就可以用同一个函数处理多种信号。显然，这也是一个回调函数，不是被\BookCode{main}函数调用，而是被系统所调用。



当某个信号的处理函数被调用时，内核自动将当前信号加入进程的信号屏蔽字，当信号处理函数返回时自动恢复原来的信号屏蔽字，这样就保证了在处理某个信号时，如果这种信号再次产生，那么它会被阻塞到当前处理结束为止。如果在调用信号处理函数时，除了当前信号被自动屏蔽之外，还希望自动屏蔽另外一些信号，则用\BookCode{sa\_\BookSplit{}mask}字段说明这些需要额外屏蔽的信号，当信号处理函数返回时自动恢复原来的信号屏蔽字。



\BookCode{sa\_\BookSplit{}flags}字段包含一些选项，本章的代码都把\BookCode{sa\_\BookSplit{}flags}设为0，\BookCode{sa\_\BookSplit{}sigaction}是实时信号的处理函数，本章不详细解释这两个字段，有兴趣的读者参考\hyperlink{db.bi01.bibli.apue}{[APUE2e]}。


\subsection{pause}
\BookAnchor{db.ch33s04.id2885627}%

\BookSource{numbers=none}{code/ch33s04/03.pv}{5}{3}


\BookCode{pause}函数使调用进程挂起直到有信号递达。如果信号的处理动作是终止进程，则进程终止，\BookCode{pause}函数没有机会返回；如果信号的处理动作是忽略，则进程继续处于挂起状态，\BookCode{pause}不返回；如果信号的处理动作是捕捉，则调用了信号处理函数之后\BookCode{pause}返回-1，\BookCode{errno}设置为\BookCode{EINTR}，所以\BookCode{pause}只有出错的返回值（想想以前还学过什么函数只有出错返回值？）。错误码\BookCode{EINTR}表示“被信号中断”。



下面我们用\BookCode{alarm}和\BookCode{pause}实现\BookCode{sleep(\BookSplit{}3)\BookSplit{}}函数，称为\BookCode{mysleep}。


\BookExampleTitle{mysleep}\BookAnchor{db.ch33s04.id2885725}

\BookSource{numbers=left,xleftmargin=19pt}{code/ch33s04/04.pv}{9}{36}

\begin{enumerate}

\item \BookCode{main}函数调用\BookCode{mysleep}函数，后者调用\BookCode{sigaction}注册了\BookCode{SIGALRM}信号的处理函数\BookCode{sig\_\BookSplit{}alrm}。

\item 调用\BookCode{alarm(\BookSplit{}nsecs)\BookSplit{}}设定闹钟。

\item 调用\BookCode{pause}等待，内核切换到别的进程运行。

\item \BookCode{nsecs}秒之后，闹钟超时，内核发\BookCode{SIGALRM}给这个进程。

\item 从内核态返回这个进程的用户态之前处理未决信号，发现有\BookCode{SIGALRM}信号，其处理函数是\BookCode{sig\_\BookSplit{}alrm}。

\item 切换到用户态执行\BookCode{sig\_\BookSplit{}alrm}函数，进入\BookCode{sig\_\BookSplit{}alrm}函数时\BookCode{SIGALRM}信号被自动屏蔽，从\BookCode{sig\_\BookSplit{}alrm}函数返回时\BookCode{SIGALRM}信号自动解除屏蔽。然后自动执行系统调用\BookCode{sigreturn}再次进入内核，再返回用户态继续执行进程的主控制流程（\BookCode{main}函数调用的\BookCode{mysleep}函数）。

\item \BookCode{pause}函数返回-1，然后调用\BookCode{alarm(\BookSplit{}0)\BookSplit{}}取消闹钟，调用\BookCode{sigaction}恢复\BookCode{SIGALRM}信号以前的处理动作。
\end{enumerate}


以下问题留给读者思考：



1、信号处理函数\BookCode{sig\_\BookSplit{}alrm}什么都没干，为什么还要注册它作为\BookCode{SIGALRM}的处理函数？不注册信号处理函数可以吗？



2、为什么在\BookCode{mysleep}函数返回前要恢复\BookCode{SIGALRM}信号原来的\BookCode{sigaction}？



3、\BookCode{mysleep}函数的返回值表示什么含义？什么情况下返回非0值？。


\subsection{可重入函数}
\BookAnchor{db.ch33s04.id2885983}%


当捕捉到信号时，不论进程的主控制流程当前执行到哪儿，都会先跳到信号处理函数中执行，从信号处理函数返回后再继续执行主控制流程。信号处理函数是一个单独的控制流程，因为它和主控制流程是异步的，二者不存在调用和被调用的关系，并且使用不同的堆栈空间。引入了信号处理函数使得一个进程具有多个控制流程，如果这些控制流程访问相同的全局资源（全局变量、硬件资源等），就有可能出现冲突，如下面的例子所示。


\BookFigure{不可重入函数}{\BookImage{images/signal.reentrancy.png}{0.909}}{\BookAnchor{db.ch33s04.id2886010}}


\BookCode{main}函数调用\BookCode{insert}函数向一个链表\BookCode{head}中插入节点\BookCode{node1}，插入操作分为两步，刚做完第一步的时候，因为硬件中断使进程切换到内核，再次回用户态之前检查到有信号待处理，于是切换到\BookCode{sighandler}函数，\BookCode{sighandler}也调用\BookCode{insert}函数向同一个链表\BookCode{head}中插入节点\BookCode{node2}，插入操作的两步都做完之后从\BookCode{sighandler}返回内核态，再次回到用户态就从\BookCode{main}函数调用的\BookCode{insert}函数中继续往下执行，先前做第一步之后被打断，现在继续做完第二步。结果是，\BookCode{main}函数和\BookCode{sighandler}先后向链表中插入两个节点，而最后只有一个节点真正插入链表中了。



像上例这样，\BookCode{insert}函数被不同的控制流程调用，有可能在第一次调用还没返回时就再次进入该函数，这称为重入，\BookCode{insert}函数访问一个全局链表，有可能因为重入而造成错乱，像这样的函数称为不可重入函数，反之，如果一个函数只访问自己的局部变量或参数，则称为可重入（\BookCode{Reentrant}）函数。想一下，为什么两个不同的控制流程调用同一个函数，访问它的同一个局部变量或参数就不会造成错乱？



如果一个函数符合以下条件之一则是不可重入的：


\begin{itemize}

\item 调用了\BookCode{malloc}或\BookCode{free}，因为\BookCode{malloc}也是用全局链表来管理堆的。

\item 调用了标准I/\BookSplit{}O库函数。标准I/\BookSplit{}O库的很多实现都以不可重入的方式使用全局数据结构。
\end{itemize}


SUS规定有些系统函数必须以线程安全的方式实现，这里就不列了，请参考\hyperlink{db.bi01.bibli.apue}{[APUE2e]}。


\subsection{sig\_atomic\_t类型与volatile限定符}
\BookAnchor{db.ch33s04.id2886197}%


在上面的例子中，\BookCode{main}和\BookCode{sighandler}都调用\BookCode{insert}函数则有可能出现链表的错乱，其根本原因在于，对全局链表的插入操作要分两步完成，不是一个原子操作，假如这两步操作必定会一起做完，中间不可能被打断，就不会出现错乱了。下一节线程会讲到如何保证一个代码段以原子操作完成。



现在想一下，如果对全局数据的访问只有一行代码，是不是原子操作呢？比如，\BookCode{main}和\BookCode{sighandler}都对一个全局变量赋值，会不会出现错乱呢？比如下面的程序：


\BookSource{numbers=none}{code/ch33s04/05.pv}{8}{6}


带调试信息编译，然后带源代码反汇编：


\BookTerminal{code/ch33s04/06.txt}{4}{2}


其中main函数的指令中有：


\BookTerminal{code/ch33s04/07.txt}{7}{5}


虽然C代码只有一行，但是在32位机上对一个64位的\BookCode{long long}变量赋值需要两条指令完成，因此不是原子操作。同样地，读取这个变量到寄存器需要两个32位寄存器才放得下，也需要两条指令，不是原子操作。请读者设想一种时序，\BookCode{main}和\BookCode{sighandler}都对这个变量\BookCode{a}赋值，最后变量\BookCode{a}的值发生错乱。



如果上述程序在64位机上编译执行，则有可能用一条指令完成赋值，因而是原子操作。如果\BookCode{a}是32位的\BookCode{int}变量，在32位机上赋值是原子操作，在16位机上就不是。如果在程序中需要使用一个变量，要保证对它的读写都是原子操作，应该采用什么类型呢？为了解决这些平台相关的问题，C标准定义了一个类型\BookCode{sig\_\BookSplit{}atomic\_\BookSplit{}t}，在不同平台的C语言库中取不同的类型，例如在32位机上定义\BookCode{sig\_\BookSplit{}atomic\_\BookSplit{}t}为\BookCode{int}类型。



在使用\BookCode{sig\_\BookSplit{}atomic\_\BookSplit{}t}类型的变量时，还需要注意另一个问题。看如下的例子：


\BookSource{numbers=none}{code/ch33s04/08.pv}{9}{10}


为了简洁，这里只写了一个代码框架来说明问题。在\BookCode{main}函数中首先要注册某个信号的处理函数\BookCode{sighandler}，然后在一个\BookCode{while}死循环中等待信号发生，如果有信号递达则执行\BookCode{sighandler}，在\BookCode{sighandler}中将\BookCode{a}改为1，这样再次回到\BookCode{main}函数时就可以退出\BookCode{while}循环，执行后续处理。用上面的方法编译和反汇编这个程序，在\BookCode{main}函数的指令中有：


\BookTerminal{code/ch33s04/09.txt}{7}{5}


将全局变量\BookCode{a}从内存读到\BookCode{eax}寄存器，对\BookCode{eax}和\BookCode{eax}做AND运算，若结果为0则跳回循环开头，再次从内存读变量\BookCode{a}的值，可见这三条指令等价于C代码的\BookCode{while(\BookSplit{}!a)\BookSplit{};\BookSplit{}}循环。如果在编译时加了优化选项，例如：


\BookTerminal{code/ch33s04/10.txt}{4}{2}


则\BookCode{main}函数的指令中有：


\BookTerminal{code/ch33s04/11.txt}{6}{4}


第一条指令将全局变量\BookCode{a}的内存单元直接和0比较，如果相等，则第二条指令成了一个死循环，注意，这是一个真正的死循环：即使\BookCode{sighandler}将\BookCode{a}改为1，只要没有影响Zero标志位，回到\BookCode{main}函数后仍然死在第二条指令上，因为不会再次从内存读取变量\BookCode{a}的值。



是编译器优化得有错误吗？不是的。设想一下，如果程序只有单一的执行流程，只要当前执行流程没有改变\BookCode{a}的值，\BookCode{a}的值就没有理由会变，不需要反复从内存读取，因此上面的两条指令和\BookCode{while(\BookSplit{}!a)\BookSplit{};\BookSplit{}}循环是等价的，并且优化之后省去了每次循环读内存的操作，效率非常高。所以不能说编译器做错了，只能说\emph{编译器无法识别程序中存在多个执行流程}。之所以程序中存在多个执行流程，是因为调用了特定平台上的特定库函数，比如\BookCode{sigaction}、\BookCode{pthread\_\BookSplit{}create}，这些不是C语言本身的规范，不归编译器管，程序员应该自己处理这些问题。C语言提供了\BookCode{volatile}限定符，如果将上述变量定义为\BookCode{volatile sig\_\BookSplit{}atomic\_\BookSplit{}t a=\BookSplit{}0;\BookSplit{}}那么即使指定了优化选项，编译器也不会优化掉对变量a内存单元的读写。



对于程序中存在多个执行流程访问同一全局变量的情况，\BookCode{volatile}限定符是必要的，此外，虽然程序只有单一的执行流程，但是变量属于以下情况之一的，也需要\BookCode{volatile}限定：


\begin{itemize}

\item 变量的内存单元中的数据不需要写操作就可以自己发生变化，每次读上来的值都可能不一样

\item 即使多次向变量的内存单元中写数据，只写不读，也并不是在做无用功，而是有特殊意义的
\end{itemize}


什么样的内存单元会具有这样的特性呢？肯定不是普通的内存，而是映射到内存地址空间的硬件寄存器，例如串口的接收寄存器属于上述第一种情况，而发送寄存器属于上述第二种情况。



\emph{\BookCode{sig\_\BookSplit{}atomic\_\BookSplit{}t}类型的变量应该总是加上\BookCode{volatile}限定符}，因为要使用\BookCode{sig\_\BookSplit{}atomic\_\BookSplit{}t}类型的理由也正是要加\BookCode{volatile}限定符的理由。


\subsection{竞态条件与sigsuspend函数}
\BookAnchor{db.ch33s04.id2886686}%


现在重新审视\hyperlink{db.ch33s04.id2885725}{例 33.\BookSplit{}2 “mysleep”}，设想这样的时序：


\begin{enumerate}

\item 注册\BookCode{SIGALRM}信号的处理函数。

\item 调用\BookCode{alarm(\BookSplit{}nsecs)\BookSplit{}}设定闹钟。

\item 内核调度优先级更高的进程取代当前进程执行，并且优先级更高的进程有很多个，每个都要执行很长时间

\item \BookCode{nsecs}秒钟之后闹钟超时了，内核发送\BookCode{SIGALRM}信号给这个进程，处于未决状态。

\item 优先级更高的进程执行完了，内核要调度回这个进程执行。\BookCode{SIGALRM}信号递达，执行处理函数\BookCode{sig\_\BookSplit{}alrm}之后再次进入内核。

\item 返回这个进程的主控制流程，\BookCode{alarm(\BookSplit{}nsecs)\BookSplit{}}返回，调用\BookCode{pause(\BookSplit{})\BookSplit{}}挂起等待。

\item 可是\BookCode{SIGALRM}信号已经处理完了，还等待什么呢？
\end{enumerate}


出现这个问题的根本原因是系统运行的时序（Timing）\BookAnchor{db.ch33s04.id2886805}并不像我们写程序时所设想的那样。虽然\BookCode{alarm(\BookSplit{}nsecs)\BookSplit{}}紧接着的下一行就是\BookCode{pause(\BookSplit{})\BookSplit{}}，但是无法保证\BookCode{pause(\BookSplit{})\BookSplit{}}一定会在调用\BookCode{alarm(\BookSplit{}nsecs)\BookSplit{}}之后的\BookCode{nsecs}秒之内被调用。由于异步事件在任何时候都有可能发生（这里的异步事件指出现更高优先级的进程），如果我们写程序时考虑不周密，就可能由于时序问题而导致错误，这叫做竞态条件（Race Condition）\BookAnchor{db.ch33s04.id2886848}。



如何解决上述问题呢？读者可能会想到，在调用\BookCode{pause}之前屏蔽\BookCode{SIGALRM}信号使它不能提前递达就可以了。看看以下方法可行吗？


\begin{enumerate}

\item 屏蔽\BookCode{SIGALRM}信号;

\item \BookCode{alarm(\BookSplit{}nsecs)\BookSplit{};\BookSplit{}}

\item 解除对\BookCode{SIGALRM}信号的屏蔽;

\item \BookCode{pause(\BookSplit{})\BookSplit{};\BookSplit{}}
\end{enumerate}


从解除信号屏蔽到调用\BookCode{pause}之间存在间隙，\BookCode{SIGALRM}仍有可能在这个间隙递达。要消除这个间隙，我们把解除屏蔽移到\BookCode{pause}后面可以吗？


\begin{enumerate}

\item 屏蔽\BookCode{SIGALRM}信号;

\item \BookCode{alarm(\BookSplit{}nsecs)\BookSplit{};\BookSplit{}}

\item \BookCode{pause(\BookSplit{})\BookSplit{};\BookSplit{}}

\item 解除对\BookCode{SIGALRM}信号的屏蔽;
\end{enumerate}


这样更不行了，还没有解除屏蔽就调用\BookCode{pause}，\BookCode{pause}根本不可能等到\BookCode{SIGALRM}信号。要是“解除信号屏蔽”和“挂起等待信号”这两步能合并成一个原子操作就好了，这正是\BookCode{sigsuspend}函数的功能。\BookCode{sigsuspend}包含了\BookCode{pause}的挂起等待功能，同时解决了竞态条件的问题，在对时序要求严格的场合下都应该调用\BookCode{sigsuspend}而不是\BookCode{pause}。


\BookSource{numbers=none}{code/ch33s04/12.pv}{5}{3}


和\BookCode{pause}一样，\BookCode{sigsuspend}没有成功返回值，只有执行了一个信号处理函数之后\BookCode{sigsuspend}才返回，返回值为-1，\BookCode{errno}设置为\BookCode{EINTR}。



调用\BookCode{sigsuspend}时，进程的信号屏蔽字由\BookCode{sigmask}参数指定，可以通过指定\BookCode{sigmask}来临时解除对某个信号的屏蔽，然后挂起等待，当\BookCode{sigsuspend}返回时，进程的信号屏蔽字恢复为原来的值，如果原来对该信号是屏蔽的，从\BookCode{sigsuspend}返回后仍然是屏蔽的。



以下用\BookCode{sigsuspend}重新实现\BookCode{mysleep}函数：


\BookSource{numbers=left,xleftmargin=19pt}{code/ch33s04/13.pv}{9}{32}


如果在调用\BookCode{mysleep}函数时\BookCode{SIGALRM}信号没有屏蔽：


\begin{enumerate}

\item 调用\BookCode{sigprocmask(\BookSplit{}SIG\_\BookSplit{}BLOCK,\BookSplit{} \&newmask,\BookSplit{} \&oldmask)\BookSplit{};\BookSplit{}}时屏蔽\BookCode{SIGALRM}。

\item 调用\BookCode{sigsuspend(\BookSplit{}\&suspmask)\BookSplit{};\BookSplit{}}时解除对\BookCode{SIGALRM}的屏蔽，然后挂起等待待。

\item \BookCode{SIGALRM}递达后\BookCode{suspend}返回，自动恢复原来的屏蔽字，也就是再次屏蔽\BookCode{SIGALRM}。

\item 调用\BookCode{sigprocmask(\BookSplit{}SIG\_\BookSplit{}SETMASK,\BookSplit{} \&oldmask,\BookSplit{} NULL)\BookSplit{};\BookSplit{}}时再次解除对\BookCode{SIGALRM}的屏蔽。
\end{enumerate}

\subsection{关于SIGCHLD信号}
\BookAnchor{db.ch33s04.id2887260}%


进程一章讲过用\BookCode{wait}和\BookCode{waitpid}函数清理僵尸进程，父进程可以阻塞等待子进程结束，也可以非阻塞地查询是否有子进程结束等待清理（也就是轮询的方式）。采用第一种方式，父进程阻塞了就不能处理自己的工作了；采用第二种方式，父进程在处理自己的工作的同时还要记得时不时地轮询一下，程序实现复杂。



其实，子进程在终止时会给父进程发\BookCode{SIGCHLD}信号，该信号的默认处理动作是忽略，父进程可以自定义\BookCode{SIGCHLD}信号的处理函数，这样父进程只需专心处理自己的工作，不必关心子进程了，子进程终止时会通知父进程，父进程在信号处理函数中调用\BookCode{wait}清理子进程即可。



请编写一个程序完成以下功能：父进程\BookCode{fork}出子进程，子进程调用\BookCode{exit(\BookSplit{}2)\BookSplit{}}终止，父进程自定义\BookCode{SIGCHLD}信号的处理函数，在其中调用\BookCode{wait}获得子进程的退出状态并打印。



事实上，由于UNIX的历史原因，要想不产生僵尸进程还有另外一种办法：父进程调用\BookCode{sigaction}将\BookCode{SIGCHLD}的处理动作置为\BookCode{SIG\_\BookSplit{}IGN}，这样\BookCode{fork}出来的子进程在终止时会自动清理掉，不会产生僵尸进程，也不会通知父进程。系统默认的忽略动作和用户用\BookCode{sigaction}函数自定义的忽略通常是没有区别的，但这是一个特例。此方法对于Linux可用，但不保证在其它UNIX系统上都可用。请编写程序验证这样做不会产生僵尸进程。

\BookNav{\hyperlink{db.ch33s03}{上一页}\enspace\textperiodcentered\enspace \hyperlink{db.ch33}{上一级}\enspace\textperiodcentered\enspace \hyperlink{db.index}{起始页}\enspace\textperiodcentered\enspace \hyperlink{db.ch34}{下一页}}
